\chapter{Forcing Input and Time Control}\label{ch:input}

The trajectories are driven by the model's stored velocity fields, which the engine reads on demand from disk
and places on a time axis. This chapter describes the reader layer and its lazy access to the forcing, the
requirement that the selected variables share a single time cadence --- the most common configuration error
on FESOM output --- the decoding of the time axis, the resolution of the integration window, the handling of
time-averaged forcing, the velocity handed to each step, and the cache through which several processes on a
node can share one copy of each snapshot. The in-memory read-ahead that serves the reconstruction of
\cref{ch:temporal} is part of this layer.

\section{The reader layer}\label{sec:input-reader}

A reader is selected by \conf{mesh.type}: \conf{fesom2}, \conf{fesom14}, \conf{fesomc}, \conf{nemo} or
\conf{healpix}. Each
provides the mesh, opens the forcing, and reads one snapshot at a time; the NEMO reader builds its mesh with
the engine's generic builder for structured grids, which turns the wet cell centres of a logically
rectangular grid into nodes and every four adjacent wet centres into a quadrilateral element. The forcing is opened \emph{lazily}: the dataset is mapped through \code{xarray}
without loading the field arrays, and a snapshot is materialized only when a step requires it.
\begin{description}[font=\normalfont\bfseries,style=nextline,leftmargin=1.5em]
  \item[FESOM2] The mesh comes from a mesh-diagnostics file (three naming conventions are recognised), an
    ASCII mesh directory or a Zarr store. For the multi-year layout --- one file per variable per year,
    \code{<var>.fesom.<year>.nc} --- the reader assembles exactly the variables the configuration names
    (velocities, surface height, sampled scalars, vertical diffusivity) and concatenates the years along time
    with \code{dask}, so a run may span decades without the files being merged on disk and naming
    \conf{unod} does not pull in fields that merely share a prefix. A directory in any other layout is
    assembled from the first matching file per velocity component and the other readable files on the mesh's
    dimensions.
  \item[FESOM~1.4] The mesh comes from the model's ASCII mesh directory or from the cell polygons of
    level-wise data (\cref{sec:mesh-fesom14}). The forcing, \conf{data.ncinput} (the files of
    \conf{ncinput\_v} and \conf{ncinput\_w} simply join it), is a file, a directory (its \code{*.nc} files in
    name order, hidden files skipped) or a glob pattern, for instance over the member directory of a CMIP6
    tree. Every file's header is read, and the files are grouped by the variables they hold, never by their
    names. Only the variables the run \emph{consumes} are opened: $u$ and $v$ (\conf{mode.var\_u},
    \conf{var\_v}, default \code{uo}, \code{vo}), in three dimensions $\Wvel$ (\code{wo}), the sea-surface
    height (\code{zos}) when a file holds it, every \conf{mode.sample\_vars} entry and, with beaching enabled,
    \conf{beaching.dry\_mask\_var}; a required variable that no file holds is refused, and the message lists
    what the files do hold. Two data forms are recognised from the dimensions: level-wise files
    $(\text{time},\text{depth},\text{node})$, such as the published CMIP6 output, and the raw writer's
    vectors over the three-dimensional nodes $(\text{time},\text{nodes\_3d})$, which are mapped onto the
    levels through the index table of the ASCII mesh. Model restart files, the mean files of an older model
    branch whose frame and time axis cannot be checked, and files whose layouts, levels or node coordinates
    disagree are refused. Every record of every variable is compared with the static wet mask of the mesh
    on every level before anything is changed; only then are $u$, $v$ and $\Wvel$ set to zero below the bed
    and sampled scalars continued downward from the column's last wet value, so that no fill value reaches a
    particle. The velocities come back on the levels, in node layout, with the placement \code{"level"}
    (\cref{sec:vert-uvdepth}); a two-dimensional run reads one level (\conf{mode.sigma\_layer}), the depth
    mean (\conf{mode.sigma\_mean}) or the surface. Without \code{zos} the surface height is zero. The
    reader resolves the time each record represents itself (\cref{sec:input-timeavg}). Refused before the
    output file exists: the \conf{em} and \conf{m1} diffusion schemes, since FESOM~1.4 output holds no
    diffusivity fields, and \conf{beaching.use\_eta} without a sea-surface height; the offline
    reconstruction of $\Wvel$ refuses FESOM~1.4 meshes, whose output holds \code{wo} itself.
  \item[FESOM-C] One NetCDF file carries the mesh and the forcing; separate files for $v$ and $w$ may be
    merged in.
  \item[HEALPix] The mesh is derived from the data store itself: its nodes are the wet pixel centres
    (\conf{mesh.healpix\_geometry: dual}, the default) or, with \conf{pixel}, the pixels are the elements
    and their corners the nodes (two-dimensional only); wet means finite velocity at the first record on
    the tracked level. Plain Zarr and NetCDF stores and
    reference stores (\code{.parq}, \code{.json}) are supported; the sea-surface height may come from a
    separate store with an identical time axis (\conf{data.ncinput\_eta}), and the vertical levels from
    \conf{mesh.levels} (\cref{sec:vert-coord}).
  \item[NEMO] The grid-T file is the input; the grid-U and grid-V files are found by substituting the grid
    token in its name and must share its time axis exactly. The first version is two-dimensional and requires
    a grid aligned with east and north.
\end{description}
A reader may also provide a hook \fn{prepare}, which checks the opened data against the mesh. It runs at every
entry: right after the data are opened, before any output file exists, and again at the entry of each
realization, after the \conf{max\_levels} cut, so it always sees the levels in use. The FESOM~1.4 hook only
checks --- the horizontal sizes, the full level axis, the node coordinates and the wet mask of the first record
of every variable --- and changes nothing. \code{kelp info} opens the data as a run does, the hook included.

Every reader that reads a vertical velocity converts it to the engine's positive-down convention
(\cref{sec:vert-wsign}; FESOM~1.4 writes \code{wo} positive up), and with \conf{velocity\_float32} stores the velocity fields in single precision to
halve their memory, the interpolation arithmetic staying in double precision.

\section{One cadence at a time}\label{sec:input-cadence}

FESOM output stores variables at several temporal cadences, set per variable in the model's configuration,
that nonetheless share a time-units string. The variables selected for a run must share a single cadence
(\cref{fig:cadence}): Lagrangian advection typically uses the daily node velocities and vertical velocity
together with the daily sea-surface height and any daily scalars, and must not mix them with monthly
element velocities. On the multi-year path the reader merges the selected variables onto a common time axis
with an exact join; a set that cannot be aligned raises an error with the hint that a mixed cadence is the
likely cause, the merged axis is checked to be strictly increasing, and a gap between years produces a
warning. On the other layouts every file attached to the velocity file (a separate $v$ or $w$, the
remaining files of a directory) is matched to the velocity records by time stamp: a file on another axis is
refused, and a finer one that contains every velocity stamp is thinned to them with a warning. A
mixed-cadence selection is easy to write and is caught only when the forcing is opened, not when the
configuration is read. The FESOM~1.4 reader applies the rule to the records it resolves: every consumed
variable must have exactly the records of $u$ (\cref{sec:input-timeavg}).

\begin{figure}[t]
\centering
\begin{tikzpicture}[scale=1.0]
  \draw[-{Stealth[length=5pt]},gray] (-0.2,0)--(8.4,0) node[right,font=\footnotesize,black]{time (one year)};
  \foreach \x in {0,0.25,...,8.0} \draw[lagblue,line width=0.6pt] (\x,1.55)--(\x,1.85);
  \node[font=\footnotesize,left,lagblue!70!black] at (-0.1,1.7) {daily};
  \node[font=\footnotesize,right] at (8.1,1.7) {\code{unod},\code{vnod},\code{w},\code{temp},\code{salt},\code{ssh}};
  \foreach \x in {0,0.667,...,8.0} \draw[lagelem,line width=0.9pt] (\x,0.85)--(\x,1.15);
  \node[font=\footnotesize,left,lagelem!70!black] at (-0.1,1.0) {monthly};
  \node[font=\footnotesize,right] at (8.1,1.0) {\code{u},\code{v} (element), MLD};
  \draw[laggreen,line width=1.1pt] (0,0.15)--(0,0.45);
  \node[font=\footnotesize,left,laggreen!60!black] at (-0.1,0.3) {annual};
  \node[font=\footnotesize,right] at (8.1,0.3) {\code{N2},\code{bolus}};
\end{tikzpicture}
\caption{An example of the cadences of one FESOM2 output configuration: the variables share one time-units
string but are sampled at different rates. A run must select variables of a single cadence; Lagrangian
advection uses the daily node-based fields (top). Mixing cadences --- for example the daily
\code{unod}/\code{vnod} with the monthly element \code{u}/\code{v} --- cannot be aligned on one time axis and
is rejected when the forcing is opened (a finer attached file that contains every velocity stamp is thinned
with a warning instead).}
\label{fig:cadence}
\end{figure}

\section{The time axis}\label{sec:input-axis}

The snapshot times are converted to seconds from the first record. A datetime axis and a calendar-aware
(\code{cftime}) axis are converted directly; a numeric axis with units of the form ``\emph{unit} since
\emph{epoch}'' is decoded with its calendar; a numeric axis with a bare unit (days, hours, minutes, seconds
and their abbreviations) is scaled by that unit. An index axis $0,1,\dots,N-1$ without units is \emph{not}
taken as time; any other numeric axis without units is read as seconds. The snapshot spacing is the common difference when all differences agree to $10^{-3}$, and
\conf{numerics.dt} defaults to it. In the discrete mode \conf{numerics.dt} must equal that spacing, and an
irregular axis is refused; the interpolated mode accepts both. When the axis cannot be decoded, \conf{numerics.file\_dt} supplies the
spacing and the axis becomes $t_i=i\,\Delta t_{\text{file}}$; without either, the interpolated mode of
\cref{sec:temp-modes} is refused before the output file is created.

\section{The integration window}\label{sec:input-window}

The window in time is specified flexibly. The start is an index into the forcing record
(\conf{time.start\_time\_idx}) or a calendar date (\conf{time.start\_date}), parsed in the calendar of the
axis. In the interpolated mode the run starts at that instant, whether or not it coincides with a stored
time; in the discrete mode the date must coincide with a stored time, and the run starts at that record. The length is a number of integration
steps (\conf{time.n\_steps}, the preferred form), an explicit end index, or a duration in days
(\conf{time.duration\_days}, only together with \conf{time.start\_date}, converted to
$\lfloor\text{days}\cdot86400/\dt\rfloor$ steps). When the run reaches
the end of the available forcing it either stops or wraps to the beginning, according to
\conf{on\_exhausted}, the latter permitting a periodic forcing to be cycled. Multiple independent realizations
may be requested (\conf{num\_runs}), each offset in start time by \conf{run\_interval} snapshots, which is the
basis for ensembles with the stochastic schemes of \cref{ch:diffusion}; each realization writes its own output
file.

\section{Time-averaged forcing}\label{sec:input-timeavg}

Forcing is often stored as time averages --- daily or monthly means --- and the engine must know this to place
the snapshots correctly on the time axis. The key \conf{time\_averaging} declares whether the fields are
averages (\conf{mean}) or instantaneous samples (\conf{instant}). For averaged fields, whose file timestamps
conventionally mark the \emph{end} of each averaging window, the engine shifts each timestamp to the window
\emph{midpoint}, $\tilde t_i=(t_{i-1}+t_i)/2$ and $\tilde t_0=t_0-(t_1-t_0)/2$, before it is used for temporal
reconstruction, so that a daily-mean field is treated as representative of the middle of its day. Instantaneous
fields are used unshifted. We recommend setting this key explicitly: when it is omitted the engine warns and
guesses \conf{mean} for FESOM2 and HEALPix and \conf{instant} otherwise, and the guess, though usually correct,
is not guaranteed. A start date is matched against the \emph{unshifted} file times. In the interpolated mode
the run then starts at the named instant (until 2026-10-04 a date equal to a stored stamp started half a
window earlier, at the shifted midpoint, while a date one second later started on time); a date in the last
half window, past the last midpoint, starts at the last record with a warning. In the discrete mode the run
starts at the matched record.

\paragraph{Record times resolved by the reader (FESOM~1.4).} The FESOM~1.4 reader does not hand on the
stamps as written. It decodes every file with its own epoch and calendar --- the published CMIP6 files start a
new epoch every decade, raw output one per job segment --- and resolves, once per variable, the instant each
record represents: the \emph{centre of its averaging window} for a time mean, the stamp itself for a
snapshot. The window comes from exactly one source, in this order. First, the explicit keys
\conf{data.time\_stamp} (where a stamp sits in its window: \conf{start}, \conf{centre} or \conf{end}, or
\conf{instant} for snapshots) and \conf{data.time\_window} (\conf{calendar\_month}, \conf{calendar\_year},
\conf{calendar\_day} or a fixed duration such as \code{"5 d"}), given together; they decide for every file,
and when a file also has time bounds or an output schedule a warning names what was overridden and the
first record where the two disagree. Second, the files' time bounds, as in the CMIP6 files: a January record
bounded by 0 and 31 days represents day 15.5, that is, 12:00 on the 16th. Third, the raw writer's
\code{output\_schedule} attribute, whose stamps are the window ends in seconds since the job's start date; the
reader enumerates the output events the schedule implies and requires the stamps to be exactly those events.
Otherwise the run is refused with a message naming the two keys. The key \conf{data.time\_offset}
(\code{"<$\pm$number> <s|min|h|d>"}) is added to the stamps, the bounds and the epoch of every file before any
of this; it serves raw output of a job that started at a time of day other than midnight, whose stamps all sit
early by that time of day, and for raw-writer files it must lie in $[0,\SI{24}{\hour})$. Calendar windows are
computed exactly in the files' calendar, and with a calendar window a stamp that does not sit where
\conf{data.time\_stamp} puts it --- on a boundary of the window, or at its centre --- is refused. A window is never estimated from the spacing of neighbouring stamps, every record the source says must
exist has to be present --- a gap, a missing month or file, refuses the run --- and every consumed variable
must resolve to exactly the records of $u$, the same instants and windows to within \SI{1}{\milli\second}. The
last rule is the reader's form of the single-cadence requirement of \cref{sec:input-cadence}: the default raw
output, which writes \code{zos} daily and the velocities monthly, is refused, and the message suggests leaving
\code{zos} out of the input.

The resolved instants become the dataset's time axis, and nothing downstream moves them again. For this reader
the engine gives no warning, applies no midpoint shift, refuses a set \conf{data.time\_averaging}, and writes
\code{time\_averaging = "resolved\_by\_reader"} to the output. The refusal is deliberate: the shift of this
section assumes stamps at the window ends, whereas the reader's instants are already window centres --- a
CMIP6 monthly stamp, for instance, already sits in the middle of its month --- so a second shift would move
every record by half a window. Daily means therefore sit at 12:00, and a \conf{start\_date} is matched against
these centres: a discrete-mode start at midnight matches no record and needs \conf{temporal\_interp}. Monthly
means lie a calendar month apart, an irregular spacing that the discrete mode refuses, so they need the
interpolated mode. The three keys are refused for every other mesh type, and the run summary and
\code{kelp info} print one line saying what resolved the times, with the first and the last instant.

\codeanchor{\raggedright The record times are \fn{resolve\_axis} in \srcfile{readers/fesom14/\_timeaxis.py}, called per
consumed variable by \fn{Fesom14Reader.open\_dataset} (\srcfile{readers/fesom14/reader.py}); the keys reach it
through \fn{\_time\_option\_kwargs}, which also refuses \conf{data.time\_averaging} for the mesh types in
\code{\_READER\_RESOLVES\_TIME}, and \fn{\_describe\_record\_times} prints the summary line
(\srcfile{readers/\_\_init\_\_.py}); \fn{run\_tracking} (\srcfile{tracking.py}) skips the midpoint shift and
records \code{RESOLVED\_BY\_READER}.}


\section{The velocity handed to a step}\label{sec:input-payload}

Snapshots are read through a small cache keyed by record index, which keeps the bracketing pair, the PCHIP
stencil and the read-ahead window and releases the rest (\cref{sec:temp-prefetch}). Between two snapshots at
fraction $\interpfrac$ the step receives a lazy view of the pair rather than an interpolated field: each field
is blended, $(1-\interpfrac)\,a+\interpfrac\,b$, the first time a consumer reads it and is then cached, and a
field that no consumer reads is never blended --- the two-dimensional fused kernels and the
three-dimensional barycentric RK2 kernel read none, the fused three-dimensional Euler kernel reads the
blended velocities, and the beaching hook the surface height and a dry mask. At a snapshot, to within $10^{-12}$ in $\interpfrac$, the snapshot itself is passed. With the
view come the target snapshot and the time to it (\cref{sec:temp-linear}). Large blends run as a parallel
loop that is bit-identical to the array expression; fields stored in single precision are blended serially.

\codeanchor{The readers live in \srcfile{readers/} (one module per model; the registry and
\fn{extract\_time\_info} in \srcfile{readers/\_\_init\_\_.py}); the multi-year assembly and its cadence check
are in \srcfile{readers/fesom2.py}; the structured-grid builder is \srcfile{readers/structured.py}; the FESOM~1.4 reader is the package
\srcfile{readers/fesom14/} (\fn{Fesom14Reader} in \srcfile{reader.py}, the two data forms and the validity
predicate in \srcfile{\_layout.py}), and the \fn{prepare} hook is run by \fn{\_prepare\_dataset}
(\srcfile{readers/\_\_init\_\_.py}); the
midpoint shift (\fn{\_shift\_to\_midpoints}), the start date (\fn{\_resolve\_start\_date}), the payload
(\fn{\_build\_vel\_payload}) and the cache \fn{\_VelCache} are in \srcfile{tracking.py}; the lazy view
\fn{LazyLerpField} and its blend are in \srcfile{core/io.py}.}

\section{Sharing snapshots between processes}\label{sec:input-shared}

A large release is often split into several processes on one node (\cref{sec:perf-shard}), and each would
otherwise read and hold its own copy of every snapshot. With \conf{numerics.shared\_fields\_dir} set, one
process reads and publishes each processed snapshot and the others map the same bytes read-only
(\cref{fig:shared-cache}). A process that finds a snapshot missing takes a file lock; the one that wins
builds the arrays, writes them under a temporary name and renames the directory into place, which is atomic,
so a sibling never sees a partial snapshot. The writer keeps a window of three snapshots on either side of
the current one and removes the rest.

The snapshots live under a namespace derived from everything that determines their content: the
configuration keys that change what a snapshot read returns (mesh, mode, the input paths, the velocity
precision, the diffusivity and dry-mask variables), the size and modification time of the mesh and data files,
the engine version, and a \emph{generation}. The generation is the key
\conf{numerics.shared\_fields\_generation} when it is set; otherwise it is taken from the environment variable
\code{KELP\_SHARED\_FIELDS\_GENERATION}, and failing that from the batch job and job step. Shards of one job therefore
share a namespace, while another job, or a rewritten data file, never reuses a snapshot; without a generation
the run falls back to local reads. Each process registers itself on start and deregisters when it ends, and a
new cache removes namespaces whose processes have all ended. Every error of the shared layer falls back to
local reads with one warning, read-ahead is disabled, and the trajectories are identical with and without the
cache. Placed on the parallel file system and shared within one node, the cache costs almost nothing; shared
across nodes it slows every step (\cref{sec:perf-io}). Node-local memory (\code{/dev/shm}) has not been
measured.

\begin{figure}[t]
\centering
\begin{tikzpicture}[node distance=4mm and 8mm]
  \node[flowbox,text width=22mm] (p1) {\footnotesize process 1\\(writer)};
  \node[flowbox,text width=22mm,below=of p1] (p2) {\footnotesize process 2};
  \node[flowbox,text width=22mm,below=of p2] (p3) {\footnotesize process $S$};
  \node[flowbox,fill=laggreen!10,draw=laggreen,text width=47mm,right=26mm of p2,minimum height=26mm] (store)
    {\footnotesize\textbf{namespace}\\[2pt] \code{t00000041/}\\ \code{t00000042/}\\ \code{t00000043.tmp} $\to$ rename\\[2pt] lock, users, window $\pm3$};
  \node[flowbox,text width=22mm,left=16mm of p2] (data) {\footnotesize forcing files};
  \draw[flowarrow] (data)--(p1.west) node[midway,above,font=\footnotesize,sloped]{read};
  \draw[flowarrow] (p1.east)--(store.north west) node[midway,above,font=\footnotesize,sloped]{publish};
  \draw[flowarrow,densely dashed] (store.west)--(p2.east) node[midway,above,font=\footnotesize]{map};
  \draw[flowarrow,densely dashed] (store.south west)--(p3.east) node[midway,below,font=\footnotesize,sloped]{map};
\end{tikzpicture}
\caption{The node-shared snapshot cache. The process that wins the lock reads a snapshot from the forcing
files, writes it under a temporary name and renames it into place; the other processes map the published
arrays read-only. The namespace is derived from the configuration, the input files, the engine version and
the job, so no other run reuses it.}
\label{fig:shared-cache}
\end{figure}

\codeanchor{\fn{SharedFieldCache} in \srcfile{core/shared\_fields.py}; the namespace is
\fn{\_shared\_fields\_tag} with \fn{\_path\_fingerprint} and \fn{\_shared\_fields\_generation}
(\srcfile{tracking.py}).}
